% Part: normal-modal-logic
% Chapter: frame-definability
% Section: equivalence-S5

\documentclass[../../../include/open-logic-section]{subfiles}

\begin{document}

\olfileid{nml}{frd}{es5}

\olsection{علاقات التكافؤ و\Log{S5}}

للمنطق الموجهي~$\Log{S5}$ وصف بالأطر الكلية: هو مجموعة
ال!!{formula}s الصالحة في كل إطار يصل فيه كل عالم إلى كل عالم،
ويدخل في ذلك وصوله إلى نفسه. وعلى هذا يُقرأ~$\Box$ ضرورة
و$\Diamond$ إمكانية؛ فصدق~$\Box !A$ هو صدق~$!A$ في \emph{كل}
عالم، وصدق~$\Diamond !A$ هو صدق~$!A$ في \emph{بعض} العوالم.
لكن~$\Log{S5}$ يقبل وصفًا آخر مكافئًا: ال!!{formula}s الصالحة
في جميع الأطر الانعكاسية المتناظرة المتعدية، أي الأطر التي علاقاتها
\emph{علاقات تكافؤ}.

\begin{defn}
  تكون العلاقة الثنائية~$\OLId{latin-looped}{R}$ على~$\OLId{latin-looped}{W}$ \emph{علاقة تكافؤ} إذا وفقط إذا
  كانت انعكاسية ومتناظرة ومتعدية. وتكون العلاقة~$\OLId{latin-looped}{R}$ على~$\OLId{latin-looped}{W}$
  \emph{كلية} إذا وفقط إذا كان~$Ruv$ لكل~$\OLId{latin-ordinary}{u},\OLId{latin-ordinary}{v} \in \OLId{latin-looped}{W}$.
\end{defn}

ومن تمييز~\Ax{T} و\Ax{B} و\Ax{4} للانعكاسية والتناظر والتعدي
نعلم أن أطر التكافؤ هي بالضبط الأطر التي تصلح فيها هذه
ال!!{formula}s الثلاث. وليس هذا السبيل الوحيد؛ إذ يمكن الجمع
بين !!{formula}s أخرى للحصول على الصنف نفسه، لأن شروط التكافؤ
تساوي في مؤداها تركيبات أخرى من خصائص~$\OLId{latin-looped}{R}$ المألوفة.

\begin{prop}\ollabel{prop:equivalences}
  الشروط الآتية متكافئة:
  \begin{enumerate}
  \item $\OLId{latin-looped}{R}$ علاقة تكافؤ؛
  \item $\OLId{latin-looped}{R}$ انعكاسية وإقليدية؛
  \item $\OLId{latin-looped}{R}$ تسلسلية ومتناظرة وإقليدية؛
  \item $\OLId{latin-looped}{R}$ تسلسلية ومتناظرة ومتعدية.
  \end{enumerate}
\end{prop}

\begin{proof}
  تمرين.
\end{proof}

\begin{prob}
  أثبت~\olref[nml][frd][es5]{prop:equivalences} بإثبات الآتي:
  \begin{enumerate}
  \item إذا كانت~$\OLId{latin-looped}{R}$ متناظرة ومتعدية، فهي إقليدية.
  \item إذا كانت~$\OLId{latin-looped}{R}$ انعكاسية، فهي تسلسلية.
  \item إذا كانت~$\OLId{latin-looped}{R}$ انعكاسية وإقليدية، فهي متناظرة.
  \item إذا كانت~$\OLId{latin-looped}{R}$ متناظرة وإقليدية، فهي متعدية.
  \item إذا كانت~$\OLId{latin-looped}{R}$ تسلسلية ومتناظرة ومتعدية، فهي انعكاسية.
  \end{enumerate}
  اشرح لماذا يكفي هذا لإثبات تكافؤ الشروط.
\end{prob}

تمثل~\olref{prop:equivalences} النظير الدلالي لـ
\olref[prf][prs]{prop:S5}، إذ تعطي وصفًا مكافئًا للمنطق الموجهي
للأطر التي تكون~$\OLId{latin-looped}{R}$ عليها علاقة تكافؤ (وهو المنطق المسمى تقليديًا
$\Log{S5}$).

ويبقى بيان الصلة بين الكلية والتكافؤ. كل علاقة كلية علاقة تكافؤ،
وليس كل تكافؤ كليًا؛ فالصنفان مختلفان. ومع ذلك يتفقان فيما يصلح
عليهما من الصيغ: ما يصلح على جميع العلاقات الكلية هو بعينه
ال!!{formula}s الصالحة على جميع علاقات التكافؤ.

\begin{prop}
  لتكن~$\OLId{latin-looped}{R}$ علاقة تكافؤ، وعرّف، لكل~$\OLId{latin-ordinary}{w} \in \OLId{latin-looped}{W}$، \emph{صنف تكافؤ}~$\OLId{latin-ordinary}{w}$
  بأنه المجموعة~$[\OLId{latin-ordinary}{w}] = \{\OLId{latin-ordinary}{w}'\in \OLId{latin-looped}{W} : Rww'\}$. عندئذ:
  \begin{enumerate}
  \item $\OLId{latin-ordinary}{w} \in [\OLId{latin-ordinary}{w}]$؛
  \item $\OLId{latin-looped}{R}$ كلية على كل صنف تكافؤ~$[\OLId{latin-ordinary}{w}]$؛
  \item مجموعة أصناف التكافؤ تقسّم~$\OLId{latin-looped}{W}$ إلى مجموعات جزئية متنافية
    اثنتين اثنتين ومستنفدة لها مجتمعة.
  \end{enumerate}
\end{prop}

\begin{prop}\ollabel{prop:S5=univ}
  تكون !!^a{formula}~$!A$ صالحة في جميع الأطر
  $\mModel{F} =\tuple{\OLId{latin-looped}{W},\OLId{latin-looped}{R}}$ التي تكون~$\OLId{latin-looped}{R}$ فيها علاقة تكافؤ، إذا وفقط
  إذا كانت صالحة في جميع الأطر~$\mModel{F} =\tuple{\OLId{latin-looped}{W},\OLId{latin-looped}{R}}$ التي تكون
  $\OLId{latin-looped}{R}$ فيها كلية. ومن ثم فإن منطق الأطر الكلية ليس إلا~$\Log{S5}$.
\end{prop}

\begin{proof}
  نبدأ بالاتجاه المباشر: العلاقة الكلية~$\OLId{latin-looped}{R}$ على~$\OLId{latin-looped}{W}$ علاقة تكافؤ.
  ومن ثم، إذا كانت~$!A$ صالحة في جميع الأطر التي تكون~$\OLId{latin-looped}{R}$ فيها علاقة
  تكافؤ، فهي صالحة في جميع الأطر الكلية. وللاتجاه الآخر نستدل
  بالمعاكس النقيض: افترض أن~$!B$ !!a{formula} تفشل في عالم~$\OLId{latin-ordinary}{w}$ من
  نموذج~$\mModel{M} = \tuple{\OLId{latin-looped}{W}, \OLId{latin-looped}{R}, \OLId{latin-looped}{V}}$ مبني على إطار~$\tuple{\OLId{latin-looped}{W},\OLId{latin-looped}{R}}$
  تكون~$\OLId{latin-looped}{R}$ عليه علاقة تكافؤ على~$\OLId{latin-looped}{W}$. إذن~$\mSat/{\OLId{latin-looped}{M}}{!B}[\OLId{latin-ordinary}{w}]$. عرّف
  نموذجًا~$\mModel{M}' = \tuple{\OLId{latin-looped}{W}', \OLId{latin-looped}{R}', \OLId{latin-looped}{V}'}$ كما يأتي:
  \begin{enumerate}
  \item $\OLId{latin-looped}{W}' = [\OLId{latin-ordinary}{w}]$؛
  \item $\OLId{latin-looped}{R}'$ كلية على~$\OLId{latin-looped}{W}'$؛
  \item $\OLId{latin-looped}{V}'(\OLId{latin-ordinary}{p}) = \OLId{latin-looped}{V}(\OLId{latin-ordinary}{p}) \cap \OLId{latin-looped}{W}'$.
  \end{enumerate}
  (إذن تمثل المنطقة المظللة في~\olref{fig:partition} مجموعة
  العوالم~$\OLId{latin-looped}{W}'$ في~$\mModel{M}'$.) ويسهل أن نرى أن~$\OLId{latin-looped}{R}$ و$\OLId{latin-looped}{R}'$ تتفقان
  على~$\OLId{latin-looped}{W}'$. ويمكن عندئذ أن نثبت بالاستقراء على ال!!{formula}s أنه لكل
  $\OLId{latin-ordinary}{w}' \in \OLId{latin-looped}{W}'$: يصدق~$\mSat{\OLId{latin-looped}{M}'}{!A}[\OLId{latin-ordinary}{w}']$ إذا وفقط إذا كان
  $\mSat{\OLId{latin-looped}{M}}{!A}[\OLId{latin-ordinary}{w}']$، لكل~$!A$ (وهذا ذو معنى لأن~$\OLId{latin-looped}{W}' \subseteq \OLId{latin-looped}{W}$).
  فحُفظ المثال المضاد عند الاقتصار على صنف التكافؤ؛ إذ بخاصة~$\mSat/{\OLId{latin-looped}{M}'}{!B}[\OLId{latin-ordinary}{w}]$، فتفشل~$!B$ في نموذج مبني على إطار كلي.
\end{proof}

\begin{figure}[t]
  \centering
  \begin{tikzpicture}[node distance=2cm, auto, thick]
    \clip [rounded corners] (0,0) -- (6,0) -- (6,4) -- (0,4) --  cycle;
    \begin{scope}
      \clip (0,0) -- (2,0)  .. controls (1.5,1.5) and (4.5,1.5)
      .. (4,4) -- (0,4) -- (0,0) -- cycle;
      \filldraw[fill=gray!40] (0,4) -- (0,2) .. controls (1.5,1.5)
      and (4.5,1.5) .. (4.5,0) -- (6,0) -- (6,4) -- (0,4) --  cycle;
    \end{scope}
    \draw [name path=line1] (2,0) .. controls (1.5,1.5) and (4.5,1.5) .. (4,4);
    \draw [name path=line2] (0,2)  .. controls (1.5,1.5) and (4.5,1.5) .. (4.5,0);
    \path [name intersections={of=line1 and line2}];
    \draw [rounded corners, use as bounding box, very thick] (0,0)
    -- (6,0) -- (6,4) -- (0,4) --  cycle;
    \node at (2,3) {$[\OLId{latin-ordinary}{w}]$} ;
    \node at (1,1) {$[\OLId{latin-ordinary}{u}]$} ;
    \node at (3,0.75) {$[\OLId{latin-ordinary}{v}]$} ;
    \node at (4.75,2) {$[\OLId{latin-ordinary}{z}]$} ;
  \end{tikzpicture}
  \caption{تقسيم~$\OLId{latin-looped}{W}$ إلى أصناف تكافؤ.}
  \ollabel{fig:partition}
\end{figure}

\end{document}
